Evaluation proceeds in three stages, each tied to a different facet of a neural computer. \textbf{Perceptual checks} come first: we compare each rollout against synchronized reference clips on held-out episodes and track generic video quality measures (e.g., \FVD, \LPIPS, \SSIM, optical-flow smoothness) without committing to a single scalar score. Surviving variants are then assessed on a \textbf{text gate} for CLI scenes, where we inspect whether rendered traces remain readable and semantically faithful to the underlying commands and outputs. Finally, a \textbf{control gate} evaluates GUI behavior by replaying predicted actions over ground-truth interfaces and inspecting pointer trajectories, click behavior, and task completion.

Instead of focusing on one metric, we use multiple evaluation approaches. For CLI, this includes text readability and optional automatic probes such as drift between rendered and reference buffers. For GUI, it combines trajectory-level quantities (e.g., pointer displacement errors, hover timings) with outcome-based probes such as whether clicks land on the intended widgets. As models and datasets mature, we expect the exact mix of metrics to evolve; the goal of this section is therefore to define the \emph{families} of signals that matter for NCs, not to lock in a particular numeric recipe.
